\uuid{MgTk}
\titre{Vecteurs}
\auteur{bodin ; hanani ; mzari}
\organisation{exo7}
\licence{CC BY-NC-SA 4.0}
\datecreate{2020-01-20}
\niveau{L1}
\module{Géométrie}
\chapitre{Géométrie de l'espace}
\sousChapitre{Produit scalaire, produit vectoriel, déterminant}
\difficulte{1}
\theme{Produit scalaire -- Produit vectoriel -- Déterminant, vecteurs}
\exo7qcmid{lille-1.geometrie-espace.001}

\contenu{
\texte{
Pour ces questions, l'espace est muni d'un repère orthonormé direct  $(O, \vec{i}, \vec{j}, \vec{k})$.
}

\begin{question}
\qclassification{141.01}
\qtags{motcle=vecteurs}
Soit $\vec{u}(1,1,1), \vec{v}(1,-1,0)$ et $\vec{w}(0,1,1)$ trois vecteurs. Quelles sont les assertions vraies ?
\begin{answers}
    \good{$\vec{u}$ et $\vec{v}$ sont orthogonaux.}

    \bad{$\vec{v}$ et $\vec{w}$ sont colinéaires.}

    \good{$(O,\vec{u},\vec{v},\vec{w})$ est un repère de l'espace.}

    \bad{$(O,\vec{u},\vec{v},\vec{w})$ est un repère orthonormé de l'espace.}
\end{answers}
\begin{explanations}
Deux vecteurs $\vec{u}$ et $\vec{v}$ sont orthogonaux si et seulement si $ \vec{u} \cdot \vec{v}=0$.
 $(O,\vec{u},\vec{v},\vec{w})$ est un repère si et seulement si $\det (\vec{u},\vec{v},\vec{w}) \neq 0$.
\end{explanations}
\end{question}
}
